\chapter{执行摘要}

\section{项目概述}

「慧农识病」是一个面向中国小农户的\hlgreen{轻量化 AI 农事参谋工具}。农户只需用手机拍摄作物叶片或虫体的照片并上传，系统即可在秒级时间内自动识别病虫害种类、输出识别置信度，并直接给出「\hlgreen{用什么药、怎么治、什么时候防}」的具体行动建议。

本项目的核心理念是做「减法」而非「加法」：不做功能臃肿的农业综合平台，而是聚焦一个高频、刚需、高价值的单一场景——病虫害诊断与防治，打造「\hlgold{拍照即诊断、诊断即给方案、给方案即能落地}」的极简工具，让小农户「零学习成本、零专业门槛」即可使用。

\section{核心优势与差异化亮点}

\begin{enumerate}[leftmargin=1.6em]
  \item \hlgreen{极轻量化，端侧可用}：采用深度可分离卷积、知识蒸馏与 INT8 量化技术，模型体积压缩至数 MB，可在中低端智能手机上离线运行，不依赖网络、不依赖算力，契合下沉市场硬件条件。
  \item \hlgreen{全闭环设计}：打通「识别——防治——预警」完整链路，不仅告诉农户「这是什么病」，更直接给出用药类型、用量、施用时机，以及基于地理定位的当季高发病虫害预防预警。
  \item \hlgreen{落地性强}：贴合小农户生产习惯，界面极简、术语白话化，与农资终端、基层农技站、村镇组织深度绑定，解决「最后一公里」推广难题。
  \item \hlgreen{数据壁垒可累积}：随着拍照数据的积累，形成「作物—区域—季节」病虫害发生地图，反哺模型迭代与预警精度，形成正向飞轮，构筑长期竞争壁垒。
\end{enumerate}

\section{市场机会与商业价值}

我国农业正加速数字化转型，智慧农业市场规模已突破万亿元量级并保持两位数增长。与此同时，农户智能手机普及率已达 76\%，但真正将其用于农业生产管理的极少，约三分之二的受访者从未关注过农业类公众号——这既是巨大的市场空白，也意味着\hlgold{巨大的结构性机会}：一端是高达数亿的小农户群体的刚需未被满足，另一端是 AI 图像识别、轻量化模型等技术已足够成熟、成本足够低。二者结合，本项目有望率先在「病虫害 AI 诊断」这一垂直细分赛道建立用户心智与规模优势。

\section{融资需求与资金用途}

本项目拟进行\hlgold{种子轮融资 50 万元}，出让 10\% 股权（投前估值 450 万元）。资金主要投向技术研发与模型迭代（40\%）、市场推广与渠道建设（24\%）、田间数据采集与标注（16\%）、团队与运营（12\%）及储备资金（8\%）。资金到位后预计支撑约 24 个月的研发与市场投入，直至实现盈亏平衡。本轮融资规模刻意控制在较小体量，以匹配学生团队当前的执行能力与项目实际需求。

\section{核心财务指标与发展预期}

\begin{table}[H]
  \centering
  \small
  \renewcommand{\arraystretch}{1.25}
  \caption{核心财务指标预测}
  \begin{tabular}{lccc}
    \toprule
    指标 & 第 1 年 & 第 2 年 & 第 3 年 \\
    \midrule
    累计注册用户（万人） & 2 & 8 & 25 \\
    营业收入（万元） & 12 & 55 & 180 \\
    净利润（万元） & $-33$ & $-10$ & 55 \\
    盈亏平衡 & \multicolumn{3}{c}{预计第 24--26 个月实现单月盈亏平衡} \\
    \bottomrule
  \end{tabular}
\end{table}

本项目不是简单的「识别工具」，而是一个\hlgreen{以极简工具为入口、以精准防治方案为价值、以数据壁垒为护城河}的农业科技项目。我们相信，在乡村振兴与农业数字化双重政策红利下，本项目具备明确的社会价值与可观的投资回报。
